\chapter{Results and Discussion}\label{ch:results-and-discussion}

At the date of this first draft, the repository substantiates subsystem
development but does not contain a complete, traceable data package
demonstrating upright balance, disturbance rejection, or
three-degree-of-freedom planar motion. The chapter therefore reports the
defensible qualitative results and provides a fixed structure for
quantitative results. Bracketed result fields must be populated from raw
data before submission. Statements that depend on those fields are
explicitly conditional.

\section{Confirmed Development Results}\label{confirmed-development-results}

The project has established a nonconventional mechanical concept: four
Mecanum wheels share a common transverse line rather than occupying the
corners of a rectangle. Four integrated CubeMars actuators and 104 mm
nominal Mecanum wheels were selected. An ESP32 firmware implementation
partitions portable interfaces and device drivers from board-specific
tasks, acquires the ISM330DHCX, performs mounting and gyro calibration
steps, produces a complementary pitch estimate, receives CRSF commands,
updates a local display, communicates with actuators through TWAI\slash{}CAN,
monitors feedback and faults, and gates commands through supervisor
logic \cite{zephyr_repo}. These are integration milestones; none independently
proves closed-loop vehicle performance.

The electrical repository includes a KiCad project and a developing
power\slash{}control design. Component selection and schematic work are more
mature than physical PCB validation. The present evidence does not
establish a fabricated and tested custom board, final protection design,
conducted or radiated compatibility, or continuous-current capability.
Similarly, the mathematical development establishes a symbolic framework
and a rank test, but numerical predictions remain conditional on
measured geometry and mass properties.

\section{Balancing Performance}\label{balancing-performance}

The primary balance result will be evaluated over a predefined interval
beginning after arming or disturbance release. Let
e\_\ensuremath{\theta}[k]=\ensuremath{\theta}\_ref[k]\ensuremath{-}\ensuremath{\theta}[k]. The RMS error is the square root of
the mean of e\_\ensuremath{\theta}\ensuremath{^{2}} over N valid samples. Settling time is the earliest
time after which the absolute pitch error remains within a stated band
for the remainder of the evaluation interval. Overshoot is reported only
when the response and reference make that quantity well defined.
Saturation time, fall-state transitions, and excluded samples are
reported alongside error metrics.

\begin{equation}
\theta_{\mathrm{RMS}}=\sqrt{\frac{1}{N}\sum_{k=1}^{N}
\left(\theta_k-\theta_{ref,k}\right)^2}
\label{eq:11-1}
\end{equation}

\clearpage
\begin{landscape}
\par\medskip\noindent\singlespacing
\phantomsection\addcontentsline{lot}{table}{\protect\numberline{11.1}Balancing and motion-performance results}
\textbf{Table 11.1. Balancing and motion-performance results}\par

\begin{longtable}[]{@{}
  >{\raggedright\arraybackslash}p{(\columnwidth - 10\tabcolsep) * \real{0.2460}}
  >{\raggedright\arraybackslash}p{(\columnwidth - 10\tabcolsep) * \real{0.1349}}
  >{\raggedright\arraybackslash}p{(\columnwidth - 10\tabcolsep) * \real{0.1984}}
  >{\raggedright\arraybackslash}p{(\columnwidth - 10\tabcolsep) * \real{0.1111}}
  >{\raggedright\arraybackslash}p{(\columnwidth - 10\tabcolsep) * \real{0.2063}}
  >{\raggedright\arraybackslash}p{(\columnwidth - 10\tabcolsep) * \real{0.1032}}@{}}
\toprule\noalign{}
\begin{minipage}[b]{\linewidth}\raggedright
\textbf{Metric}
\end{minipage} & \begin{minipage}[b]{\linewidth}\raggedright
\textbf{Simulation}
\end{minipage} & \begin{minipage}[b]{\linewidth}\raggedright
\textbf{Hardware mean \ensuremath{\pm} SD}
\end{minipage} & \begin{minipage}[b]{\linewidth}\raggedright
\textbf{Trials}
\end{minipage} & \begin{minipage}[b]{\linewidth}\raggedright
\textbf{Acceptance criterion}
\end{minipage} & \begin{minipage}[b]{\linewidth}\raggedright
\textbf{Status}
\end{minipage} \\
\midrule\noalign{}
\endhead
\bottomrule\noalign{}
\endlastfoot
RMS pitch error (deg) & \textless INSERT\textgreater{} &
\textless INSERT\textgreater{} & \textless INSERT N\textgreater{} &
\textless AUTHOR CONFIRM\textgreater{} & Pending \\
Peak pitch excursion (deg) & \textless INSERT\textgreater{} &
\textless INSERT\textgreater{} & \textless INSERT N\textgreater{} &
\textless AUTHOR CONFIRM\textgreater{} & Pending \\
Settling time (s) & \textless INSERT\textgreater{} &
\textless INSERT\textgreater{} & \textless INSERT N\textgreater{} &
\textless AUTHOR CONFIRM\textgreater{} & Pending \\
Steady-state pitch bias (deg) & \textless INSERT\textgreater{} &
\textless INSERT\textgreater{} & \textless INSERT N\textgreater{} &
\textless AUTHOR CONFIRM\textgreater{} & Pending \\
Maximum recoverable initial angle (deg) & \textless INSERT\textgreater{}
& \textless INSERT\textgreater{} & \textless INSERT N\textgreater{} &
\textless AUTHOR CONFIRM\textgreater{} & Pending \\
Longitudinal velocity RMSE (m\slash{}s) & \textless INSERT\textgreater{} &
\textless INSERT\textgreater{} & \textless INSERT N\textgreater{} &
\textless AUTHOR CONFIRM\textgreater{} & Pending \\
Lateral velocity RMSE (m\slash{}s) & \textless INSERT\textgreater{} &
\textless INSERT\textgreater{} & \textless INSERT N\textgreater{} &
\textless AUTHOR CONFIRM\textgreater{} & Pending \\
Yaw-rate RMSE (rad\slash{}s) & \textless INSERT\textgreater{} &
\textless INSERT\textgreater{} & \textless INSERT N\textgreater{} &
\textless AUTHOR CONFIRM\textgreater{} & Pending \\
Peak current per branch (A) & \textless INSERT\textgreater{} &
\textless INSERT\textgreater{} & \textless INSERT N\textgreater{} &
Hardware limit & Pending \\
Worst control-cycle execution (\ensuremath{\mu}s) & \textless INSERT\textgreater{} &
\textless INSERT\textgreater{} & \textless INSERT N\textgreater{} &
Below period & Pending \\
\end{longtable}
\end{landscape}
\clearpage

\begin{longtable}[]{@{}
  >{\raggedright\arraybackslash}p{(\columnwidth - 0\tabcolsep) * \real{1.0000}}@{}}
\toprule\noalign{}
\begin{minipage}[b]{\linewidth}\raggedright
\textless INSERT FIGURE 11.1: Plot time (s) horizontally. Use aligned
panels for pitch reference and measured pitch (deg), pitch rate (deg\slash{}s),
collective balance command and individual actuator commands (documented
units), saturation\slash{}fault flags, and operating state. Show all trials
faintly and the mean with uncertainty; identify the metric interval and
controller configuration\textgreater{}\\
What the figure should show: Plot time (s) horizontally. Use aligned
panels for pitch reference and measured pitch (deg), pitch rate (deg\slash{}s),
collective balance command and individual actuator commands (documented
units), saturation\slash{}fault flags, and operating state. Show all trials
faintly and the mean with uncertainty; identify the metric interval and
controller configuration\\
Likely source: synchronized raw balance-test logs processed by a
version-controlled script\strut
\end{minipage} \\
\midrule\noalign{}
\endhead
\bottomrule\noalign{}
\endlastfoot
\end{longtable}

\begin{center}\singlespacing
\phantomsection\addcontentsline{lof}{figure}{\protect\numberline{11.1}Measured pitch response during upright balancing.}
\textbf{Figure 11.1. Measured pitch response during upright balancing.}
\end{center}

\begin{longtable}[]{@{}
  >{\raggedright\arraybackslash}p{(\columnwidth - 0\tabcolsep) * \real{1.0000}}@{}}
\toprule\noalign{}
\begin{minipage}[b]{\linewidth}\raggedright
\textless INSERT RESULT: RMS pitch error, maximum pitch excursion,
settling time, overshoot, steady-state error, saturation duration, and
number of successful\slash{}attempted trials from the completed balancing
experiment. State the calculation interval and
uncertainty.\textgreater{}
\end{minipage} \\
\midrule\noalign{}
\endhead
\bottomrule\noalign{}
\endlastfoot
\end{longtable}

\section{Disturbance Rejection}\label{disturbance-rejection}

Disturbance trials will be grouped by applied impulse or initial angle
rather than pooled indiscriminately. The discussion will distinguish
recovery without saturation, recovery with saturation, transition to
FALLEN, and mechanical restraint contact. A stronger conclusion---that
the robot rejects disturbances up to a specified magnitude---will be
made only after the disturbance input is measured or repeatably defined.

\begin{longtable}[]{@{}
  >{\raggedright\arraybackslash}p{(\columnwidth - 0\tabcolsep) * \real{1.0000}}@{}}
\toprule\noalign{}
\begin{minipage}[b]{\linewidth}\raggedright
\textless INSERT FIGURE 11.2: Plot time from disturbance (s) versus
pitch angle (deg), pitch rate (deg\slash{}s), applied or proxy disturbance,
four actuator commands, battery voltage (V), and operating state.
Compare at least five repetitions at each disturbance level and mark
recovery\slash{}fall classification\textgreater{}\\
What the figure should show: Plot time from disturbance (s) versus pitch
angle (deg), pitch rate (deg\slash{}s), applied or proxy disturbance, four
actuator commands, battery voltage (V), and operating state. Compare at
least five repetitions at each disturbance level and mark recovery\slash{}fall
classification\\
Likely source: instrumented pendulum\slash{}impulse test or controlled
initial-displacement logs\strut
\end{minipage} \\
\midrule\noalign{}
\endhead
\bottomrule\noalign{}
\endlastfoot
\end{longtable}

\begin{center}\singlespacing
\phantomsection\addcontentsline{lof}{figure}{\protect\numberline{11.2}Measured recovery from repeatable pitch disturbances.}
\textbf{Figure 11.2. Measured recovery from repeatable pitch disturbances.}
\end{center}

\section{Planar Motion and Mobility}\label{planar-motion-and-mobility}

Longitudinal, lateral, and yaw tests will first be interpreted
independently. Reference-to-measurement delay, steady-state gain, RMS
error, cross-axis response, pitch deviation, and actuator effort will be
reported. Pure lateral response is especially important: wheel rotation
alone is insufficient evidence because roller deformation and slip can
generate large odometric motion without equivalent chassis translation.
External pose data are therefore required for a quantitative
lateral-mobility claim.

\begin{longtable}[]{@{}
  >{\raggedright\arraybackslash}p{(\columnwidth - 0\tabcolsep) * \real{1.0000}}@{}}
\toprule\noalign{}
\begin{minipage}[b]{\linewidth}\raggedright
\textless INSERT FIGURE 11.3: Create panels for longitudinal velocity
v\_x (m\slash{}s), lateral velocity v\_y (m\slash{}s), yaw rate (rad\slash{}s), x-y
trajectory (m), pitch (deg), four wheel speeds (rad\slash{}s), and allocation
residual. Compare reference, external measurement, wheel-based estimate,
and simulation at a common sample interval\textgreater{}\\
What the figure should show: Create panels for longitudinal velocity
v\_x (m\slash{}s), lateral velocity v\_y (m\slash{}s), yaw rate (rad\slash{}s), x-y
trajectory (m), pitch (deg), four wheel speeds (rad\slash{}s), and allocation
residual. Compare reference, external measurement, wheel-based estimate,
and simulation at a common sample interval\\
Likely source: motion-capture\slash{}camera data synchronized with firmware and
simulation logs\strut
\end{minipage} \\
\midrule\noalign{}
\endhead
\bottomrule\noalign{}
\endlastfoot
\end{longtable}

\begin{center}\singlespacing
\phantomsection\addcontentsline{lof}{figure}{\protect\numberline{11.3}Measured longitudinal, lateral, and yaw responses.}
\textbf{Figure 11.3. Measured longitudinal, lateral, and yaw responses.}
\end{center}

\begin{longtable}[]{@{}
  >{\raggedright\arraybackslash}p{(\columnwidth - 0\tabcolsep) * \real{1.0000}}@{}}
\toprule\noalign{}
\begin{minipage}[b]{\linewidth}\raggedright
\textless INSERT RESULT: For separate and combined commands, report
longitudinal-velocity RMSE, lateral-velocity RMSE, yaw-rate RMSE,
cross-axis coupling, trajectory endpoint error, maximum pitch deviation,
allocation residual, and saturation fraction.\textgreater{}
\end{minipage} \\
\midrule\noalign{}
\endhead
\bottomrule\noalign{}
\endlastfoot
\end{longtable}

\section{Electrical, Thermal, Communication, and Timing Results}\label{electrical-thermal-communication-and-timing-results}

Vehicle performance depends on electrical and computational margins.
Voltage sag can reduce actuator authority, temperature can trigger
derating, stale CAN feedback can invalidate the state used by the
controller, and scheduling jitter can reduce phase margin. The final
analysis will place these signals on the same time axis as pitch and
commands. Datasheet limits will be identified separately from
project-configured limits and measured values.

\begin{longtable}[]{@{}
  >{\raggedright\arraybackslash}p{(\columnwidth - 0\tabcolsep) * \real{1.0000}}@{}}
\toprule\noalign{}
\begin{minipage}[b]{\linewidth}\raggedright
\textless INSERT FIGURE 11.4: Plot time (s) versus battery and
regulated-rail voltage (V), total and per-branch current (A), cumulative
energy (Wh), actuator\slash{}controller\slash{}battery temperature (\ensuremath{^{\circ}}C), and commanded
operating state. Include ambient temperature and limit lines from the
verified configuration\textgreater{}\\
What the figure should show: Plot time (s) versus battery and
regulated-rail voltage (V), total and per-branch current (A), cumulative
energy (Wh), actuator\slash{}controller\slash{}battery temperature (\ensuremath{^{\circ}}C), and commanded
operating state. Include ambient temperature and limit lines from the
verified configuration\\
Likely source: power analyzer, firmware telemetry, and calibrated
temperature logs\strut
\end{minipage} \\
\midrule\noalign{}
\endhead
\bottomrule\noalign{}
\endlastfoot
\end{longtable}

\begin{center}\singlespacing
\phantomsection\addcontentsline{lof}{figure}{\protect\numberline{11.4}Motor current, battery voltage, and actuator temperature during testing.}
\textbf{Figure 11.4. Motor current, battery voltage, and actuator temperature during testing.}
\end{center}

\begin{longtable}[]{@{}
  >{\raggedright\arraybackslash}p{(\columnwidth - 0\tabcolsep) * \real{1.0000}}@{}}
\toprule\noalign{}
\begin{minipage}[b]{\linewidth}\raggedright
\textless INSERT FIGURE 11.5: Plot histograms and time histories of
control-cycle period and execution time (\ensuremath{\mu}s), IMU sample age (\ensuremath{\mu}s),
motor-feedback age (\ensuremath{\mu}s), CAN utilization (\%), error frames,
retries\slash{}timeouts, and task overruns. Report median, 95th, 99th, and
worst-case values\textgreater{}\\
What the figure should show: Plot histograms and time histories of
control-cycle period and execution time (\ensuremath{\mu}s), IMU sample age (\ensuremath{\mu}s),
motor-feedback age (\ensuremath{\mu}s), CAN utilization (\%), error frames,
retries\slash{}timeouts, and task overruns. Report median, 95th, 99th, and
worst-case values\\
Likely source: instrumented firmware time stamps and CAN analyzer
capture\strut
\end{minipage} \\
\midrule\noalign{}
\endhead
\bottomrule\noalign{}
\endlastfoot
\end{longtable}

\begin{center}\singlespacing
\phantomsection\addcontentsline{lof}{figure}{\protect\numberline{11.5}Measured control-loop timing and CAN communication performance.}
\textbf{Figure 11.5. Measured control-loop timing and CAN communication performance.}
\end{center}

\section{Requirement Verification}\label{requirement-verification}

\clearpage
\begin{landscape}
\par\medskip\noindent\singlespacing
\phantomsection\addcontentsline{lot}{table}{\protect\numberline{11.2}Requirement-verification results}
\textbf{Table 11.2. Requirement-verification results}\par

\begin{longtable}[]{@{}
  >{\raggedright\arraybackslash}p{(\columnwidth - 8\tabcolsep) * \real{0.2016}}
  >{\raggedright\arraybackslash}p{(\columnwidth - 8\tabcolsep) * \real{0.2419}}
  >{\raggedright\arraybackslash}p{(\columnwidth - 8\tabcolsep) * \real{0.1694}}
  >{\raggedright\arraybackslash}p{(\columnwidth - 8\tabcolsep) * \real{0.2661}}
  >{\raggedright\arraybackslash}p{(\columnwidth - 8\tabcolsep) * \real{0.1210}}@{}}
\toprule\noalign{}
\begin{minipage}[b]{\linewidth}\raggedright
\textbf{Requirement}
\end{minipage} & \begin{minipage}[b]{\linewidth}\raggedright
\textbf{Evidence required}
\end{minipage} & \begin{minipage}[b]{\linewidth}\raggedright
\textbf{Observed result}
\end{minipage} & \begin{minipage}[b]{\linewidth}\raggedright
\textbf{Pass criterion}
\end{minipage} & \begin{minipage}[b]{\linewidth}\raggedright
\textbf{Verdict}
\end{minipage} \\
\midrule\noalign{}
\endhead
\bottomrule\noalign{}
\endlastfoot
R-F01 controlled upright balance & Free-balance logs &
\textless INSERT\textgreater{} & Approved pitch\slash{}settling limits &
Pending \\
R-F02 commanded planar motion & External pose and commands &
\textless INSERT\textgreater{} & Approved tracking limits & Pending \\
R-C01 rank-adequate allocation & Measured H and singular values &
\textless INSERT\textgreater{} & rank(H)=3 with acceptable conditioning
& Pending \\
R-E01 protected power distribution & Schematic\slash{}layout\slash{}bench test &
\textless INSERT\textgreater{} & No limit violation or fault &
Pending \\
R-S01 safe loss-of-link response & Injected timeout test &
\textless INSERT\textgreater{} & Deterministic safe transition &
Pending \\
R-S02 stale-feedback response & Injected delay\slash{}dropout &
\textless INSERT\textgreater{} & Commands inhibited within limit &
Pending \\
R-P01 timing & Timestamp distribution & \textless INSERT\textgreater{} &
Worst case below approved deadline & Pending \\
R-M01 serviceability & Inspection\slash{}replacement demonstration &
\textless INSERT\textgreater{} & Approved procedure\slash{}time & Pending \\
\end{longtable}
\end{landscape}
\clearpage

\section{Failure Cases and Limitations}\label{failure-cases-and-limitations}

The expected failure modes include incorrect IMU or wheel signs,
estimator drift or vibration sensitivity, actuator saturation, feedback
latency, CAN node misconfiguration, voltage sag, tire\slash{}roller slip,
compliance, insufficient yaw leverage, a rank-deficient roller
arrangement, and a center of mass outside the modeled range. A result is
not removed merely because it is unfavorable; the associated
configuration and evidence are retained so the limitation can guide
redesign.

\clearpage
\begin{landscape}
\par\medskip\noindent\singlespacing
\phantomsection\addcontentsline{lot}{table}{\protect\numberline{11.3}Failure and anomaly log}
\textbf{Table 11.3. Failure and anomaly log}\par

\begin{longtable}[]{@{}
  >{\raggedright\arraybackslash}p{(\columnwidth - 10\tabcolsep) * \real{0.1176}}
  >{\raggedright\arraybackslash}p{(\columnwidth - 10\tabcolsep) * \real{0.1471}}
  >{\raggedright\arraybackslash}p{(\columnwidth - 10\tabcolsep) * \real{0.1912}}
  >{\raggedright\arraybackslash}p{(\columnwidth - 10\tabcolsep) * \real{0.1471}}
  >{\raggedright\arraybackslash}p{(\columnwidth - 10\tabcolsep) * \real{0.1985}}
  >{\raggedright\arraybackslash}p{(\columnwidth - 10\tabcolsep) * \real{0.1985}}@{}}
\toprule\noalign{}
\begin{minipage}[b]{\linewidth}\raggedright
\textbf{Event\slash{}time}
\end{minipage} & \begin{minipage}[b]{\linewidth}\raggedright
\textbf{Configuration}
\end{minipage} & \begin{minipage}[b]{\linewidth}\raggedright
\textbf{Observed symptom}
\end{minipage} & \begin{minipage}[b]{\linewidth}\raggedright
\textbf{Evidence}
\end{minipage} & \begin{minipage}[b]{\linewidth}\raggedright
\textbf{Root cause}
\end{minipage} & \begin{minipage}[b]{\linewidth}\raggedright
\textbf{Corrective action\slash{}retest}
\end{minipage} \\
\midrule\noalign{}
\endhead
\bottomrule\noalign{}
\endlastfoot
\textless INSERT\textgreater{} & \textless COMMIT\slash{}BUILD\textgreater{} &
\textless INSERT\textgreater{} & \textless LOG\slash{}PHOTO\textgreater{} &
\textless CONFIRMED OR HYPOTHESIS\textgreater{} &
\textless INSERT\textgreater{} \\
\textless INSERT\textgreater{} & \textless COMMIT\slash{}BUILD\textgreater{} &
\textless INSERT\textgreater{} & \textless LOG\slash{}PHOTO\textgreater{} &
\textless CONFIRMED OR HYPOTHESIS\textgreater{} &
\textless INSERT\textgreater{} \\
\textless INSERT\textgreater{} & \textless COMMIT\slash{}BUILD\textgreater{} &
\textless INSERT\textgreater{} & \textless LOG\slash{}PHOTO\textgreater{} &
\textless CONFIRMED OR HYPOTHESIS\textgreater{} &
\textless INSERT\textgreater{} \\
\end{longtable}
\end{landscape}
\clearpage

The present evidence supports the conclusion that Zephyr is a
substantially developed mechatronic prototype with operational
subsystems and a coherent path to validation. It does not yet support
the stronger conclusions that the robot balances autonomously, achieves
holonomic motion, meets quantitative requirements, or is electrically
and thermally qualified. Those conclusions remain pending the tests and
result fields in this chapter.
